\documentclass[a4paper]{article}
\usepackage{graphics,color,eurosym,latexsym}
\usepackage{times}
\usepackage{listings}
\lstset{columns=fixed,basicstyle=\ttfamily,numbers=left,numberstyle=\tiny,stepnumber=5,breaklines=true}
%\usepackage[german]{babel}
%\pagestyle{empty}
\usepackage[utf8]{inputenc}
\usepackage[T1]{fontenc}
\usepackage{eurosym}
\usepackage[round]{natbib}
\bibliographystyle{plainnat}
\oddsidemargin=0cm
\evensidemargin=0cm
\headheight=0cm
\headsep=0cm
\newcommand{\be}{\begin{enumerate}}
\newcommand{\ee}{\end{enumerate}}
\newcommand{\bi}{\begin{itemize}}
\newcommand{\ei}{\end{itemize}}
\newcommand{\I}{\item}
\newcommand{\ty}{\texttt}
\textwidth=16cm
\textheight=23cm
\begin{document}

\title{Comments on Alfy Code}
\author{Bernhard Haubold}
\maketitle

\section{Sliding Window Analysis}
For the sliding window analysis it might be helpful to distinguish
more clearly between a \emph{window} and an \emph{interval}. A
\emph{window} has a fixed length and is slid across the query. An
\emph{interval} initially has the same length as the window, but can
then be expanded. The sliding window analysis should result in a list
of non-overlapping intervals. So what is currently called ``window'',
I will call ``interval'' here to help me think.

After storage of an interval in the \ty{else} phrase of chunk 43b, a
new interval is created that is identical to the current window
$(\ell,r)$. Shouldn't the new interval start at $r+1$?

\bibliography{/home/haubold/References/references}

\end{document}

